% Header aus der Vorlage
\documentclass[a4paper,11pt, footheight=26pt
%,twoside
]{scrreprt}
\usepackage[head=23pt]{geometry}	% head=23pt umgeht Fehlerwarnung, dafür größeres "top" in geometry
\geometry{a4paper, top=30mm, bottom=22mm,headsep=10mm, footskip=12mm
, left=20mm, right=20mm
%, inner=27mm, outer=13mm
}

\setcounter{secnumdepth}{3}	% zählt auch subsubsection
\setcounter{tocdepth}{3}	% Inhaltsverzeichnis bis in subsubsection

% Zeile 2 (,twoside) und 7 (inner=...) für eine Druckversion (doppelseitig) ent-kommentieren (Rand für Hefter)

% Input inkl. Umlaute, Silbentrennung
\usepackage[T1]{fontenc}
\usepackage[utf8]{inputenc}
\usepackage[ngerman]{babel}
\usepackage{csquotes}	% Anführungszeichen
\usepackage{eurosym}

% HTW Corporate Design: Arial (Helvetica)
\usepackage{helvet}
\renewcommand{\familydefault}{\sfdefault}

% Style-Aufhübschung
\usepackage{soul, color}	% Kapitälchen, Unterstrichen, Durchgestrichen usw. im Text
\usepackage{scrlayer-scrpage}	% Kopf-/Fußzeile
%\usepackage{titleref}
\usepackage[perpage]{footmisc}	% Fußnotenzählung Seitenweit, nicht Dokumentenweit
\renewcommand*{\thefootnote}{\fnsymbol{footnote}}	% Fußnoten-Symbole anstatt Zahlen
\renewcommand*{\titlepagestyle}{empty} % Keine Seitennummer auf Titelseite

% Mathe usw.
\usepackage{amssymb}
\usepackage[fleqn]{amsmath}	% fleqn: align-Umgebung rechtsbündig
\usepackage{xcolor}
\usepackage{esint}	% Schönere Integrale, \oiint vorhanden
\everymath=\expandafter{\the\everymath\displaystyle}	% Mathe Inhalte werden weniger verkleinert
\usepackage{wasysym}	% mehr Symbole, bspw \lightning
% Auch arcus-Hyperbolicus-Funktionen
\DeclareMathOperator{\arccot}{arccot}
\DeclareMathOperator{\arccosh}{arccosh}
\DeclareMathOperator{\arcsinh}{arcsinh}
\DeclareMathOperator{\arctanh}{arctanh}
\DeclareMathOperator{\arccoth}{arccoth} 
% Mathe in Anführungszeichen:
\newsavebox{\mathbox}\newsavebox{\mathquote}
\makeatletter
\newcommand{\mq}[1]{% \mathquotes{<stuff>}
  \savebox{\mathquote}{\text{"}}% Save quotes
  \savebox{\mathbox}{$\displaystyle #1$}% Save <stuff>
  \raisebox{\dimexpr\ht\mathbox-\ht\mathquote\relax}{"}#1\raisebox{\dimexpr\ht\mathbox-\ht\mathquote\relax}{''}
}
\makeatother

% tikz usw.
\usepackage{tikz}
\usepackage{pgfplots}
\pgfplotsset{compat=1.11}	% Umgeht Fehlermeldung
\usetikzlibrary{graphs}
%\usetikzlibrary{through}	% ???
\usetikzlibrary{arrows}
\usetikzlibrary{arrows.meta}	% Pfeile verändern / vergrößern: \draw[-{>[scale=1.5]}] (-3,5) -> (-3,3);
\usetikzlibrary{automata,positioning} % Zeilenumbruch im Node node[align=center] {Text\\nächste Zeile} automata für Graphen
\usetikzlibrary{matrix}
\usetikzlibrary{patterns}	% Schraffierte Füllung
\tikzstyle{reverseclip}=[insert path={	% Inverser Clip \clip
	(current page.north east) --
	(current page.south east) --
	(current page.south west) --
	(current page.north west) --
	(current page.north east)}
% Nutzen: 
%\begin{tikzpicture}[remember picture]
%\begin{scope}
%\begin{pgfinterruptboundingbox}
%\draw [clip] DIE FLÄCHE, IN DER OBJEKT NICHT ERSCHEINEN SOLL [reverseclip];
%\end{pgfinterruptboundingbox}
%\draw DAS OBJEKT;
%\end{scope}
%\end{tikzpicture}
]	% Achtung: dafür muss doppelt kompliert werden!
\usepackage{graphpap}	% Grid für Graphen
\tikzset{every state/.style={inner sep=2pt, minimum size=2em}}

% Tabular
\usepackage{longtable}	% Große Tabellen über mehrere Seiten
\usepackage{multirow}	% Multirow/-column: \multirow{2[Anzahl der Zeilen]}{*[Format]}{Test[Inhalt]} oder \multicolumn{7[Anzahl der Reihen]}{|c|[Format]}{Test2[Inhalt]}
\renewcommand{\arraystretch}{1.3} % Tabellenlinien nicht zu dicht
\usepackage{colortbl}
\arrayrulecolor{gray}	% heller Tabellenlinien

% Nützliches
\usepackage{verbatim}	% u.a. zum auskommentieren via \begin{comment} \end{comment}
\usepackage{tabto}	% Tabs: /tab zum nächsten Tab oder /tabto{.5 \CurrentLineWidth} zur Stelle in der Linie
\NumTabs{6}	% Anzahl von Tabs pro Zeile zum springen
\usepackage{listings} % Source-Code mit Tabs
\usepackage{lstautogobble} 
\usepackage{enumitem}	% Anpassung der enumerates
\setlist[enumerate,1]{label=\arabic*.)}	% global andere Enum-Items
\renewcommand{\labelitemiii}{$\scriptscriptstyle ^\blacklozenge$} % global andere 3. Item-Aufzählungszeichen
\usepackage{letltxmacro} % neue Definiton von Grundbefehlen
% Nutzen:
%\LetLtxMacro{\oldemph}{\emph}
%\renewcommand{\emph}[1]{\oldemph{#1}}

% Einrichtung von lst
\lstset{
basicstyle=\ttfamily, 
mathescape=true, 
%escapeinside=^^, 
autogobble, 
tabsize=2,
basicstyle=\footnotesize\sffamily\color{black},
frame=single,
rulecolor=\color{lightgray},
numbers=left,
numbersep=5pt,
numberstyle=\tiny\color{gray},
commentstyle=\color{gray},
keywordstyle=\color{green},
stringstyle=\color{orange},
morecomment=[l][\color{magenta}]{\#}
%showspaces=false,
showstringspaces=false,
breaklines=true,
literate=%
    {Ö}{{\"O}}1
    {Ä}{{\"A}}1
    {Ü}{{\"U}}1
    {ß}{{\ss}}1
    {ü}{{\"u}}1
    {ä}{{\"a}}1
    {ö}{{\"o}}1
    {~}{{\textasciitilde}}1
}
\usepackage{scrhack} % Fehler umgehen
\def\ContinueLineNumber{\lstset{firstnumber=last}} % vor lstlisting. Zum wechsel zum nicht-kontinuierlichen muss wieder \StartLineAt1 eingegeben werden
\def\StartLineAt#1{\lstset{firstnumber=#1}} % vor lstlisting \StartLineAt30 eingeben, um bei Zeile 30 zu starten
\let\numberLineAt\StartLineAt

% BibTeX
\usepackage[backend=bibtex, bibencoding=ascii]{biblatex}	% BibTeX
\usepackage{makeidx}
%\makeglossary
%\makeindex

% Grafiken
\usepackage{graphicx}
\usepackage{epstopdf}	% eps-Vektorgrafiken einfügen

% pdf-Setup
\usepackage{pdfpages}
\usepackage[bookmarks,%
bookmarksopen=false,% Klappt die Bookmarks in Acrobat aus
colorlinks=true,%
linkcolor=black,%
citecolor=red,%
urlcolor=green,%
]{hyperref}

% Titel, Autor usw. werden vor dem Anfang des Dokuments in einem Rutsch definiert…
\newcommand{\DTitel}[1]{\newcommand{\Dokumententitel}{#1}}
\newcommand{\DUntertitel}[1]{\newcommand{\Dokumentenuntertitel}{#1}}
\newcommand{\DAutor}[1]{\newcommand{\Dokumentenautor}{#1}}
\newcommand{\DNotiz}[1]{\newcommand{\Dokumentennotiz}{#1}}
\newcommand{\DSign}[1]{\newcommand{\Dokumentensignatur}{#1}}
\DSign{\footnotesize{\textcolor{darkgray}{Mitschrift von\\ \Dokumentenautor}}}
\newcommand{\Autorformat}[1]{\textcolor{darkgray}{Mitschrift von #1}}
% … Deswegen folgendes erst Nach Dokumentenbeginn ausführen:
\AtBeginDocument{
	\hypersetup{
		pdfauthor={\Dokumentenautor},
		pdftitle={HTW Dresden | \Dokumententitel - \Dokumentenuntertitel},
	}
	\automark[section]{section}
	\automark*[subsection]{subsection}
	\pagestyle{scrheadings}
	\ihead{\includegraphics[height=1.7em]{\workingdir LaTeX_master/HTW-Logo.eps}}
	\ohead{\Dokumententitel}
	\cfoot{\pagemark}
	\ofoot{\Dokumentensignatur}
	% Titelseite
	\title{\includegraphics[width=0.35\textwidth]{\workingdir LaTeX_master/HTW-Logo.eps}\\\vspace{0.5em}
	\Huge\textbf{\Dokumententitel} \\\vspace*{0,5cm}
	\Large \Dokumentenuntertitel \\\vspace*{4cm}}
	\author{\Autorformat{\Dokumentenautor} \vspace*{1cm}\\\Dokumentennotiz}
}

%% EIGENE BEFEHLE

% Arbeitsordner (in Abhängigkeit vom Master)
\newcommand{\workingdir}{../}

% Farbdefinitionen
\definecolor{red}{RGB}{180,0,0}
\definecolor{green}{RGB}{75,160,0}
\definecolor{blue}{RGB}{0,75,200}
\definecolor{orange}{RGB}{255,128,0}
\definecolor{yellow}{RGB}{255,245,0}
\definecolor{purple}{RGB}{75,0,160}
\definecolor{cyan}{RGB}{0,160,160}
\definecolor{brown}{RGB}{120,60,10}

\definecolor{itteny}{RGB}{244,229,0}
\definecolor{ittenyo}{RGB}{253,198,11}
\definecolor{itteno}{RGB}{241,142,28}
\definecolor{ittenor}{RGB}{234,98,31}
\definecolor{ittenr}{RGB}{227,35,34}
\definecolor{ittenrp}{RGB}{196,3,125}
\definecolor{ittenp}{RGB}{109,57,139}
\definecolor{ittenpb}{RGB}{68,78,153}
\definecolor{ittenb}{RGB}{42,113,176}
\definecolor{ittenbg}{RGB}{6,150,187}
\definecolor{itteng}{RGB}{0,142,91}
\definecolor{ittengy}{RGB}{140,187,38}

% Textfarbe ändern
\newcommand{\tred}[1]{\textcolor{red}{#1}}
\newcommand{\tgreen}[1]{\textcolor{green}{#1}}
\newcommand{\tblue}[1]{\textcolor{blue}{#1}}
\newcommand{\torange}[1]{\textcolor{orange}{#1}}
\newcommand{\tyellow}[1]{\textcolor{yellow}{#1}}
\newcommand{\tpurple}[1]{\textcolor{purple}{#1}}
\newcommand{\tcyan}[1]{\textcolor{cyan}{#1}}
\newcommand{\tbrown}[1]{\textcolor{brown}{#1}}

% Umstellen der Tabellen Definition
\newcommand{\mpb}[1][.3]{\begin{minipage}{#1\textwidth}\vspace*{3pt}}
\newcommand{\mpe}{\vspace*{3pt}\end{minipage}}

\newcommand{\resultul}[1]{\underline{\underline{#1}}}
\newcommand{\parskp}{$ $\\}	% new line after paragraph
\newcommand{\corr}{\;\widehat{=}\;}
\newcommand{\mdeg}{^{\circ}}

\newcommand{\nok}[2]{\begin{pmatrix}#1\\#2\end{pmatrix}}	% n über k
\newcommand{\mtr}[1]{\begin{pmatrix}#1\end{pmatrix}}	% Matrix
\newcommand{\dtr}[1]{\begin{vmatrix}#1\end{vmatrix}}	% Determinante (Betragsmatrix)
\renewcommand{\vec}[1]{\underline{#1}}	% Vektorschreibweise
\newcommand{\imptnt}[1]{\colorbox{red!30}{#1}}	% Wichtiges

% Definition von Titel, Autor usw.
\DTitel{Mathematik I}
\DUntertitel{Vorlesungsskript}
\DAutor{Falk Jonatan Strube}
\DNotiz{Vorlesung von Herrn Meinhold}

\begin{document}

\pagenumbering{Roman}

\maketitle
\newpage
\tableofcontents
\newpage

\pagenumbering{arabic}

\part{Elementare Grundlagen}

\section{Aussagen und Grundzüge der Logik}
%\input{HTW_Mathe1_Skript_sec1}

\section{Mengen}\label{sec:Mengen}
%\input{HTW_Mathe1_Skript_sec2}

\section{Zahlen}
\subsection{Gruppen, Ringe, Körper}
\begin{itemize}
\item Gegeben sei eine Menge $M$ und eine zweistellige Operation $\circ$ (d.h. Abb. von $M\times M$ in $M$)\\
Bezeichnung: $(M,\circ)$, analog $(M, \circ, *)$
\item Die Operation $\circ$ heißt \emph{kommutativ}, wenn $a \circ b= b \circ a $ für alle $a, b \in M$.
\item Die Operation $\circ$ heißt \emph{assoziativ}, wenn $(a \circ b) \circ c =a \circ (b \circ c)$ für alle  $a, b, c \in M$.
\end{itemize}

\paragraph{Def. 1:}\parskp
$(M,\circ)$ heißt \emph{Gruppe}, wenn gilt:
\begin{enumerate}
\item Die Operation $\circ$ ist assoziativ
\item Es gibt genau ein \emph{neutrales Element} $e\in M$ mit $a\circ e = e\circ a = a$ (für alle $a \in M$)
\item Es gibt zu jedem $a \in M$ genau ein \emph{inverses Element} $a^{-1}$ mit $a\circ a^{-1}=a^{-1}\circ a = e$
\item Eine Gruppe heißt \emph{ABELsch}, wenn zusätzlich folgendes gilt:\\
$\circ$ ist kommutativ
\end{enumerate}

\paragraph{Def. 2:}\parskp
$(M,\oplus,*)$ heißt \emph{Ring}, wenn gilt:
\begin{enumerate}
\item $(M, \oplus)$ ist eine ABELsche Gruppe.
\item Die Operation $*$ ist assoziativ.
\item Es gelten für beliebige $a,b,c\in M$:\\
$a*(b\oplus c)=(a*b)\oplus (a*c)\\
(a \oplus b) *c = (a*c) \oplus (b*c)$ \qquad (Distributivgesetze)
\item Ein Ring heiß \emph{kommutativer Ring}, wenn gilt:\\
$*$ ist kommutativ
\end{enumerate}

\paragraph{Def. 3:}\parskp
$(M, \oplus, *)$ heißt \emph{Körper}, wenn gilt:
\begin{enumerate}
\item $(M,\oplus, *)$ ist ein Ring \\
(mit dem neutralen Element $E_0$ für die Operation $\oplus$)
\item $(M\setminus \{E_0\}, *)$ ist eine ABELsche Gruppe\\
(mit dem neutralen Element $E_1$ für die Operation $*$)
\end{enumerate}

\subsection{Zahlentheorie}
\begin{itemize}
\item Eine natürliche Zahl $p>1$, die nurch durch $1$ und sich selbst teilbar ist heißt \emph{Primzahl}.
\item Jede natürliche Zahl $n>1$ ist entweder eine Primzahl, oder sie lässt sich als Produkt von Primzahlen schreiben.\\
Diese sogenannte \emph{Primfaktorzerlegung} ist bis auf die Reihenfolge der Faktoren eindeutig.
\end{itemize}

\paragraph{Def. 4:}\parskp
Zwei natürliche zahlen aus $\mathbb{N}^*$ heißen \emph{teilerfremd}, wenn sie außer $1$ keine gemeinsamen teiler besitzen.
\begin{itemize}
\item Es sei $a \in \mathbb{Z}$ und $m\in\mathbb{N}^*$. Dann gibt es eine eindeutige Darstellung der Gestalt $\boxed{a=q\cdot m + r}$ mit $\boxed{0\leq r < m}$ und $q \in \mathbb{Z}$.\\
Bezeichnung: $m$ … \emph{Modul} \qquad $r$ … (kleinste nichtnegative) \emph{Rest modulo m} ($r \equiv mod (a,m)$)
\item Zur Erinnerung: $a$ und $b$ seien ganze Zahlen, $m\in \mathbb{n}^*$, dann $a\equiv b (mod\; m)$ [$a$ kongruent $b\; modulo\; m$]\\
$\Leftrightarrow$ $a$ und $b$ haben den gleicher Rest $modulo\; m$\\
$\Leftrightarrow$ $a-b$ ist durch $m$ teilbar (d.h. $\exists k \in \mathbb{Z} \quad a-b = k\cdot m$)
\paragraph{Satz 1:} \parskp
Es sei $a \equiv b (mod\; m)$, $c\equiv d (mod\; m)$, dann gilt: $a+c \equiv b+d (mod \;m)$ und $a\cdot c \equiv b \cdot d (mod\; m)$ (d.h. in Summen und Produktenn darf jede Zahl durch einen beliebigen Vertreter der gleichen Restklasse ersetzt werden).
\subparagraph{Bsp. 1:} 
\begin{enumerate}[label=\alph*)]
\item $307+598 \equiv 1+(-2)\equiv -1 \equiv 5 (mod \; 6)$
\item $307\cdot 598 \equiv 1 \cdot (-2) \equiv -2 \equiv 4 (mod \; 6)$
\item $598^6 \equiv (-2)^6\equiv 64 \equiv 4 (mod \; 6)$
\end{enumerate}
\item Man wählt aus jeder Restklasse den kleinsten nichtnegativen Vertreter \\
$\curvearrowright$ Menge von Resten $modulo \; m$: $\mathbb{Z}_m:= \{0,1,...,m-1\}$ \\
$\curvearrowright$ „modulare Arithmetik“: Operation $\oplus$ und $\odot$ für Zahlen aus $\mathbb{Z}_m$ erklärbar, in dem für das Ergebnis jeweils der kleinste nichtnegative Rest $modulo \; m$ gewählt wird (vgl. Satz 1)\\
z.B. $\mathbb{Z}_7 = \{ 0,1,..., 6\}$, \quad $5\oplus 4=2$, da $5+4\equiv 9\equiv 2 (mod \; 7)$ \quad $5\odot 6=2$, da $5\cdot 6\equiv 30 \equiv 2 (mod \;7)$\smallskip\\
Falls keine Verwechselung zu befürchten ist, wird die übliche Schreibweise $+$ und $\cdot$ anstelle von $\oplus$ und $\odot$ verwendet.
\end{itemize}

\paragraph{Def. 5:} \parskp
Wenn es zu $c \in \mathbb{Z}_m$ eine Zahl $d\in \mathbb{Z}_m$ gibt, mit $c\cdot d \equiv 1 (mod \; m)$ (bzw. $c\odot d \equiv 1$), so heißt $d$ die \emph{(multiplikative) modulare Inverse} zu $c$ in $\mathbb{Z}_m$.\\
Bezeichnung: $d=c^{-1}$
\subparagraph{Bsp. 2:} \parskp
$c=3 \in \mathbb{Z}_7$, wegen $3 \cdot 5 \equiv 1 (mod \; 7)$ ist (in $\mathbb{Z}_7$) $3^{-1}=5$.
\subparagraph{Satz 2:} Zu $a \in \mathbb{Z}_m, a \not = 0$, gibt es genau dann eine modulare Inverse in $\mathbb{Z}_m$, wenn $a$ und $m$ teilerfremd sind ($ggT(a,m)=1$).
\subparagraph{Satz 3:} Es sei $p$ eine Primzahl. Dann ist $(\mathbb{Z}_m, \oplus, \odot)$ ein Körper.\\
Bemerkung: Falls $m$ keine Primzahl ist, so ist $(\mathbb{Z}_m, \oplus, \odot)$ ein kommutativer Ring.
\subparagraph*{EUKLIDischer Algorithmus}
\begin{itemize}
\item Verfahren zur Ermittlung des größten gemeinsamen Teilers $t$ zweier positiver natürlicher Zahlen, $t=ggT(a,b)$.
\item In erweiterter Form bietet der Algorithmus eine Möglichkeit zur Bestimmung  der modularen Inversen von $a$ zum Modul $m$ (mit $a<m$ und $a,m$ teilerfremd).
\end{itemize}

\subparagraph{Satz 4:} (EUKLIDischer Algorithmus)\\
Es seien $a,b \in \mathbb{N}^*, a>b$. Man bildet die endliche Folge \\
$r_0:= b, \; r_1=mod \;(a,b), \;r_2 = mod \; (r_0,r_1),...,\;r_n = mod \; (r_{n-2}, r_{n-1})$, Abbruch falls $r_n=0$.\\
In diesem Fall gist $\boxed{ggT(a,b)=r_{n-1}}$ (letzter nicht verschwindender Rest).\\
Bezeichnung: $j$-te Division ... $\boxed{r_{j-2}:r_{j-1}= q_j\text{ Rest } r_j}$ ($j=1,...,n$) (dabei $r_1:=a$).

\subparagraph{Satz 5:}  (erweiterter EUKLIDischer Algorithmus)\\
Zusätzlich zur Folge $(r_n)$ aus Satz 4 bilde man die Folgen \\
$x_0=0,\; x_1=1, \; x_2 = x_0 - q_2x_1, ...,\; x_j=x_{j-2}-q_jx_{j-1}\quad (j\leq n-1)$ und\\
$y_0=1, \; y_1=-q_1, \; y_2=y_0- q_2 y_1, ... ,\; y_j=y_{j-2}-q_jy_{j-1}\quad (j \leq n-1)$\\
Dann gilt für alle $j=0, ..., \; n-1$: $\boxed{r_j=x_j\cdot a + y_j \cdot b}$\\
Insbesondere gilt $\boxed{ggT(a,b) = x_{n-1}\cdot a + y_{n-1}\cdot b}$
\subparagraph{Diskussion:}
\begin{enumerate}
\item Der Sinn der erweiterten EUKLIDischen Algorithmus besteht darin, in jedem Schrit den \emph{Divisionsrest $r$ als linearkombination von $a$ und $b$ mit ganzzahligen Koeffizienten $x$ und $y$} darzustellen: $r=x\cdot a + y\cdot b$\\
Der Mechanismus wird am besten im Rechenschema des nachfolgenden Bsp. 4 deutlich.
\item  Sind $c$ und $m$ teilerfremd, $1\leq c < m$, d.h. $ggT(m,c)=1$, so erhält man mit dem erweiterten EUKLIDischen Algorithmus ($a=m, b=c$) eine Darstellung in der Form $\boxed{1=x\cdot m+ y\cdot c}$.\\
$\curvearrowright y\cdot c \equiv 1 (mod\; m)$ und damit $c^{-1}\equiv y (mod\; m)$ (für die modulare Inverse muss eventuell noch der in $\mathbb{Z}_m$ liegende, zu $y$ kongruente, Wert gebildet werden!).
\end{enumerate}

\subparagraph{Bsp. 3:} \parskp
Man ermittle den größten gemeinsamen Teiler $t$ sowie das kleinste gemeinsame Vielfache $v$ der Zahlen $132$ und $84$.
\begin{itemize}
\item Es genügt der „einfache“ Algorithmus:
\begin{align*}
132:84 &= 1\text{ Rest } 48\\
84:48 &= 1\text{ Rest } 36\\
48:36 &= 1\text{ Rest } 12 \quad \curvearrowright t = ggT (132,84) = \resultul{12}\\
36:12 &= 3\text{ Rest } \boxed{0} \curvearrowright \text{ Ende.}\\
\end{align*}
\item $v=\frac{a \cdot b}{t}=\frac{132\cdot 84}{12}=\resultul{924}=kgV (132,84)$
\end{itemize}

\subparagraph{Bsp. 4:} \parskp
Man ermittle die modulare Inverse von $\overbrace{11}^b$ zum Modul $\overbrace{25}^a$. \medskip\\
\begin{tabular}{r l | l | r l}
				&				&						&	$11$&$=0 \cdot 25 + 1 \cdot 11$ (1)\\
$25:11 $&$= 2 \text{ Rest }  3$	&	$3=25-2 \cdot 11$	&	$3 $&$= 1 \cdot 25 - 2 \cdot 11$ (2)\\
$11: 3 $&$= 3 \text{ Rest } 2$		&	$2= 11-3 \cdot 3$	&	$2 $&$= -3 \cdot 25 + 7 \cdot 11$ (3)\\
$3:2 $&$= 1 \text{ Rest } 1$		&	$1 = 3-2$			&	$\boxed{1} $&$= 4 \cdot 25 - 9 \cdot 11$\\
$2 : 1 $&$= 2 \text{ Rest }  0$	&						& &\\
\end{tabular} \medskip\\
$\curvearrowright (-9) \cdot 11 \equiv 1 (mod \; 25)\\
\curvearrowright 11^{-1} \equiv -9 \equiv 16 (mod \; 25)$, die Inverse von $11$ in $\mathbb{Z}_{25}$ ist $16$.\\
Zu den Schritten:
\begin{itemize}
\item [(1)] $b=0 \cdot a + 1 \cdot b$
\item [(2)] mittleres Feld als Linearkombination
\item [(3)] ab hier Rechnung links spaltenweise durchführen, dabei Faktoren $a$ und $b$ beibehalten. 
\end{itemize}
\subsubsection*{EULERsche $\varphi$-Funktion, Satz von EULER}
\paragraph{Def. 6:} \parskp
Es sei $n \in \mathbb{N}^*$. Dann \emph{EULERsche $\varphi$-Funktion}:

$\varphi(n):=$ Anzahl der zu $n$ teilerfremden Elemente aus $\{1,2,...,n\}$.
Eigenschaften der $\varphi$-Funktion:
\begin{itemize}
\item Es sei $p$ eine Primzahl, dann ist $\boxed{\varphi(p)=p-1}$, $\boxed{\varphi(p^k)=p^{k-1}(p-1)} \; (k \in \mathbb{N}^*$
\item Falls $ggT (m, n) =1$, so gilt $\varphi (m\cdot n) = \varphi (m) \cdot \varphi (n)$.
\item Speziell: $n=p\cdot q $ ($p,q$ Primzahlen), dann $\boxed{\varphi(n)=(p-1)\cdot (q-1)}$ (1).
\end{itemize}
\subparagraph{Satz 6:} (Satz von EULER)\\
Es sei $ggT(a,n)=1$, dann gilt:\\
$\boxed{a^{\varphi(n)}\equiv 1 (mod \; n)}$ (2).

\subsubsection*{RSA-Verschlüsselung}
\begin{itemize}
\item Die Formeln (1) und (2) [siehe oberhalb] bilden die Grundlage für die sogenannte RSA-Ver\-schlüsselung (RIVES, SHAMIR, ADLEMAN - 1978)
\item Schlüsselerzeugung:
\begin{enumerate}
\item Man wählt (in der Praxis sehr große) Primzahlen $d$ und $q$.
\item $n:= p \cdot q$, $m:= \varphi (n) \overset{(1)}{=} (p-1)(q-1)$
\item $e$ wird so gewählt, dass $ggT(e,m)=1$
\item $d:= e^{-1} (mod \; m)$ (modulare Inverse)
\item $(n,e)$ … öffentlicher Schlüssel\\
$(n,d)$ … geheimer Schlüssel (geheim ist nur $d$)\\
$p$, $q$ und $m$ werden nicht mehr benötigt, bleiben aber geheim!
\end{enumerate}
\item Verschlüsselung: \\
Klartext $a$ teilerfremd zu $n$ verschlüsseln mit $e$, d.h. $b:\equiv a^e (mod \; n)$ bilden ($b$ … Geheimtext)
\item Entschlüsselung: \\
Der Empfänger und Besitzer des geheimen Schlüssels bildet $b^d (mod \; m)$ und erhält $b^d\equiv a (mod \; n)$ denn $b^d\equiv (a^e)^d \equiv a^{ed}\equiv a ^{1+k\cdot m}\equiv a\cdot \underset{\equiv 1 \text{ wegen Satz 5}}{\left( a^{\varphi(n)}\right)}^k\equiv a (mod \; n)$.
\item Praktische Durchführung vgl Übungsaufgabe 2.4
\end{itemize}

\subsection{Reelle Zahlen}

$\mathbb{R}$ … Menge der reellen Zahlen.\\
Auf $\mathbb{R}$ existiert eine algebraische Struktur und eine Ordnungsstruktur.

\subsubsection{Algebraische Struktur}
$(\mathbb{R}, + ,\cdot)$ mit den arithmetischen Operationen $+$ (Addition) und $\cdot$ (Multiplikation) ist ein Körper.
\paragraph{Def. 7:}
\begin{enumerate}[label=\alph*.)]
\item $\boxed{0! := 1, \; n!=n\cdot (n-1)!}$ mit $n\in \mathbb{N}^*$ \emph{Fakultät} (rekursive Funktion)
\item Sei $\alpha \in \mathbb{R}, k\in \mathbb{N}^*$, dann sei $\boxed{\nok{\alpha}{0}:=1, \nok{\alpha}{k}:= \frac{\alpha}{k}\nok{\alpha-1}{k-1}}$ Binominalkoeffizient $\alpha$ über $k$.\\
d.h. $\boxed{\nok{\alpha}{k}= \frac{\alpha(\alpha-1)(\alpha-2)...(\alpha-k-1)}{k!}}$
\end{enumerate}

\subparagraph{Diskussion:}
\begin{enumerate}
\item Für $k, n \in \mathbb{N}, \; 0 \leq k \leq n$ gilt $\boxed{\nok{n}{k}=\nok{n}{n-k}=\frac{n!}{k!(n-k)!}}$
\item Für $k \in \mathbb{N}, \alpha \in \mathbb{R}$ gilt $\boxed{\nok{\alpha}{k}+\nok{\alpha}{k+1}=\nok{\alpha+1}{k+1}}$
\item \emph{Binomischer Lehrsatz}: $(a+b)^k=\sum_{k=0}^n \nok{n}{k} a^{n-k} \cdot b^k=a^n+n\cdot a^{n-1} \cdot b + \nok{n}{2}a^{n-2}\cdot b^2+...+b^n$\\
\emph{Pascalsches Dreieck}:
\begin{tabular}{c c c c c c c c c c c r}
   &   &   &   &   & 1 &   &   &   &   &   & \textcolor{red}{$n=0$} \\
   &   &   &   & 1 &   & 1 &   &   &   &   & \textcolor{red}{$1$} \\
   &   &   & 1 &   & 2 &   & 1 &   &   &   & \textcolor{red}{$2$} \\
   &   & 1 &   & 3 &   & 3 &   & 1 &   &   & \textcolor{red}{$3$} \\
   & 1 &   & 4 &   & 6 &   & 4 &   & 1 &   & \textcolor{red}{$4$} \\
 1 &   & 5 &   & 10&   & 10&   & 5 &   & 1 & \textcolor{red}{$5$} \\

\end{tabular}
(Spalte $\hat{=}k$ in $\nok{n}{k}$)
\end{enumerate}

\paragraph{Stellenwertsysteme:}
\begin{itemize}
\item Es sei $k>1$ eine natürliche Zahl (die sogenannte Basis)
\item $x = (\underbrace{x_p x_{p-1}...x_1x_0}_{\text{Vorkomma}},\;\underbrace{x_{-1}x_{-2}...x_{-q}}_{\text{Nachkomma}})_b\\
:= \underbrace{x_p\cdot b^p+x_{p-1}\cdot b^{p-1}+...+x_1\cdot b^1+x_0\cdot b^0}_{\text{Vorkomma}}+\underbrace{x_{-1}\cdot b^{-1}+x_{-2}\cdot b^{-2}+...+x_{-q}\cdot b^{-q}}_{\text{Nachkomma}}$\\
heißt Darstellung zur Basis b (*).
\end{itemize}

\subparagraph{Bsp. 5:}
\begin{itemize}
\item $b=2$ … \emph{Dual- oder Binärsystem} (Ziffern $\{0,1\}$)
\item $b=3$ … Trialsystem
\item $b=10$ … \emph{Dezimalsystem}
\item $b=16$ … \emph{Hexadezimalsystem} (Ziffern $\{0,1,2,...,9,\underset{10}{A},\underset{11}{B},\underset{12}{C},\underset{13}{D},\underset{14}{E},\underset{15}{F}\}$
\end{itemize}
z.B. $(47)_{10}=(101111)_2=(1202)_3=(57)_8=(2F)_{16}$
\paragraph{Übergang von einem Ziffernsystem zu einem anderen}\parskp
z.B. $p=3,\; q=2$
\begin{align*}
x&=x_3b^3+x_2b^2+x_1b^1+x_0+x_{-1}b^{-1}+x_{-2}b^{-2}\\
&= (x_3b^2+x_2x^1+x_1)b+x_0+(x_{-1}+x_{-2}b^{-1})b^{-1}\\
&=((x_3b^1+x_2)b+x_1)b+x_0+(x_{-1}+x_{-2}b^{-1})b^{-1}
\end{align*}
Grundlage: fortgesetztes Klammern:
\begin{align*}
x&=((....((x_p b+x_{p-1})b+x_{p-2})b+...+x_2)b+x_1)b+x_0+\\
&\qquad((...(x_{-q}b^{-1}+x_{-(q-1)})b^{-1}+...+x_{-2})b^{-1}+x_{-1})b^{-1}
\end{align*} (**)
\subparagraph{Bsp. 6:} Übergang Dezimalsystem $\rightarrow$ anderes System
\begin{itemize}
\item ganzer Teil: fortgesetzte Division durch $b$ und Restabspaltung liefert $b$-Ziffern in der Reihenfolge $x_0,x_1,x_2,...$
\item gebrochener Teil: fortgesetzte Multiplikation mit $b$ und Abspaltung des ganzzahligen Anteils liefert $b$-Ziffern in der Reihenfolge $x_{-1},x_{-2},...$
\end{itemize}
z.B. Dezimalsystem $\rightarrow$ Hexadezimalsystem ($b=16$)\\
$x=435,9$\\
ganzer Teil:\\
\begin{tabular}{r c l l}
$435:16=$ & $27$ & Rest $3$ &  $\rightarrow x_0$\\
$27:16=$ & $1$ & Rest $11$ &  $\rightarrow x_1$\\
$1:16=$ & $\boxed{0}$ & Rest $1$ &  $\rightarrow x_2$\\
\end{tabular}\\
gebrochener Teil: \\
\begin{tabular}{r c l l}
$0,9 \cdot 16=$ & $0.4$ & $+14$ & $\rightarrow x_{-1}$\\
$0,4 \cdot 16=$ & $\boxed{0.4}$ & $+6$ & $\rightarrow x_{-2}$ (Periode, da gleicher „Nachkommarest“)\\
\end{tabular}\\
$\curvearrowright x=(1B3,E\overline{6})_{16}$

\subparagraph{Bsp. 7:} Übergang anderer Systeme $\rightarrow$ Dezimalsystem\\
Entweder direkte Auswertung von (*) (v.a. beim Dualsystem $\rightarrow$ Addition von 2er-Potenzen) \emph{oder} Klammern in (**) von innen nach außen berechnen (zweckmäßig HORNER-Schema).
\begin{itemize}
\item ganzer Teil: beginnend mit $x_p$\\
ABB1 
\item gebrochener Teil: beginnend mit $x_{-q}$\\
ABB 2
\end{itemize}
$x=(1E2,B8)_{16}$\\
ganzer Teil:\\
\begin{tabular}{c | c c c}
 & $1$  &  $14$  & $2$\\
 $16$  &   & $16$  & $480$\\
 \hline
  & $1$ & $30$ & $\boxed{482}$\\

\end{tabular}\\
gebrochener Teil:\\
\begin{tabular}{c | c c c}
 & $8$  &  $11$  & *\\
 $\frac{1}{16}=0,0625$  &   & $0,5$  & $\boxed{0,71875}$\\
 \hline
  & $8$ & $11,5$ & \\
\end{tabular}\\
$\curvearrowright x=\resultul{482,71875}$

\subparagraph{Bsp. 8:} Hexadezimalsystem $\leftrightarrow$ Dualsystem\\
4 Dualziffern entsprechen einer Hexadezimalziffer ($2^4=16$) $\curvearrowright$ 4er Gruppen von Dualziffern ab Komma bilden.
\begin{enumerate}[label=\alph*.)]
\item $(A8C,B2)_{16}=(1010 \; 1000 \; 1100, \; 1011 \; 1011 \; 001(0))_2$
\item $((0)110\; 1110, \; 101(0))_2 = (6E,A)_{16}$
\end{enumerate}

\subsubsection{Zahlendarstellung im Computer}

\begin{enumerate}
\item Ganze Binärzahlen in Zweierkomplementdarstellung ($n$ Bit, meist $n=8,16,32,64$)
\begin{itemize}
\item Bsp.: $n=8 \qquad (100)_{^10}=(64)_{16}$\\
\begin{tabular}{|c | c | c |c | c | c |c | c |}
\hline
\textcolor{red}{0} & 1 &1&0&0&1&0&0\\
\hline
$2^7$&$2^6$&$2^5$&$2^4$&$2^3$&$2^2$&$2^1$&$2^0$\\
\hline
\end{tabular}\\
\textcolor{red}{MSB}: most siginficant bit\\
(LSB: least significant bit)
\item Um auch negative Zahlen darstellen zu können, wird das MSB als Vorzeichen reserviert. Negative Zahlen $-a$ ($1\leq a \leq 2^{n-1}$) werden im sogenannten Zweierkomplement $\overline{a}:=2^n-a$ dargestellt ($\curvearrowright \overline{a}\geq 2^{n-1} \curvearrowright MSB =1$)
\item Nightnegtavie Zahlen $0\leq a \leq 2^{n-1}-1$ werden unverändert dargestellt ($MSB=0$)
\item Damit Darstellung ganzer Zahlen von $-2^{n-1}$ bis $2^{n-1}-1$ (Anzahl $2^n$) möglich, $n=8: -128 bis 127$.
\item Umwandlung negativer Zahlen $\rightarrow$ Zweierkomplement
\subparagraph{Bsp. 9:} $n=8$, umzuwandeln sei $-100$ (dezimal)\\
zwei Möglichkeiten:
\begin{enumerate}[label=\arabic*.)]
\item  (für die Handrechnung): $\overline{100}=\underbrace{2^8}_{=256}-100=\underbrace{156}_{(9C)_{16}}=$\begin{tabular}{|c|c c c|c c c c|}
\hline
\textcolor{red}{$1$}&$0$&$0$&$1$&$1$&$1$&$0$&$0$\\
\hline
\end{tabular}\\
\emph{Bemerkung}: Das Zweierkomplement der positiven Zahl $100$ ist die positive Zahl $156=\overline{100}$, diese wird wegen $MSB=1$ als negative Zahl $-100$ interpretiert.
\item (am schnellsten): Rechts (beim LSB) beginnend alle Ziffern bis einschließlich der ersten $1$ übernehmen (unverändert lassen), für alle höherwertigen Ziffern $Z$ das \emph{Einerkomplement} $1-z$ bilden:\\
$(100)_{10}=$\begin{tabular}{|c|c c c|c c c c|}
\hline
$0$&$1$&$1$&$0$&$0$&$1$&$0$&$0$\\
\hline
\end{tabular}\\
\begin{tabular}{|c|c c c|c c c c|}
\hline
\textcolor{red}{$1$}&\textcolor{red}{$0$}&\textcolor{red}{$0$}&\textcolor{red}{$1$}&\textcolor{red}{$1$}&\textcolor{green}{$1$}&\textcolor{green}{$0$}&\textcolor{green}{$0$}\\
\hline
\end{tabular} $=(\overline{100})_{10}=(156)_2$\\
Rückumwandlung (Zahl mit $MSB=1 \rightarrow$ neg. Zahl) analog:\\
$\overline{156}=256-156=100\rightarrow \resultul{-100}$\\
Die Subtraktion wird damit auf die Addition des Zweierkomplements zurückgeführt.
\end{enumerate}
\subparagraph{Bsp. 10:} $a=64-100=64+(-100)$
\begin{align*}
64=2^6 &= 0100 \, 0000\\
-100 &= 1001 \, 1100 \quad +\\
Summe &= \textcolor{red}{1}101\, 1100 \quad \text{\textcolor{red}{Ergebins negativ}}\\
36 &= 0010\, 0100 \quad \text{dargestellst ist aber }\overline{z}=2^n-z
\end{align*}
$\curvearrowright$ Ergebnis: $a=-36=-z$\\
Ein \emph{Überlauf} (Ergebnis $\geq 2^{n-1}$ oder $<-2^{n-1}$) ensteht in folgenden Fällen ($\rightarrow$ ERROR!)
\begin{tabular}{r | c |c |c l}
 & $a$ & $b$ & $a+b$ & \\
 \hline
 MSB & 0 & 0& \textcolor{red}{1} & ($MSB = 0$ erwartet!)\\
 \hline 
 MSB & 1 & 1& \textcolor{red}{0} & ($MSB = 1$ erwartet!)\\
\end{tabular}\\
Bemerkung: Für die Handrechnung (z.B. $2-5=:a$) kleinere zahl von größerer Subtrahieren $a=-(5-2)$, dabei genügt es für $n$ die Binärstellenzahl des Minuenden $(5)_{10}=(101)_2$ also $n=3$  zu verwenden. Es wird dabei ausschließlich mit nicht-negativen Zahlen gerechnet ($0,1,...,2^n-1$):\\
$(5-2)_{10}=((5+\textcolor{red}{2^n}-2)-\textcolor{red}{2^n})_{10}=(5+\overline{2}-\textcolor{red}{2^n})_{10}$\\
$(2)_{10}=(010)_2 \curvearrowright \overline{2}=(110)_2$\\
$(5)_{10}=(101)_2$\\
\begin{tabular}{r l}
$5$:& $101$\\
$\overline{2}$:& $110 \quad +$\\
$\textcolor{red}{1}$& $011$
\end{tabular}\\
\textcolor{red}{vordere Stelle $2^n$ ignorieren}\\
$\curvearrowright 5-3=3=(011)_2 \curvearrowright \resultul{a=-3}$
\end{itemize}
\item Gleitkommasysteme\\
$\boxed{x=v\cdot m \cdot b^e}$ dabei
\begin{itemize}
\item $v=(-1)^V$ … \emph{Vorzeichen} $\begin{cases}
V= 0 \quad \text{positive Zahl}\\
V=1 \quad \text{negative Zahl}
\end{cases}$
\item $m$ … \emph{Mantisse}, Stellenzahl $p$, die Mantisse heißt \emph{normalisiert} falls sie folgende Gestalt besitzt:\\
 $m_1, m_2, ..., m_p$ oder $0, m_1, m_2, ..., m_p$ mit $m\not = 0$. Dabei sind  $m_1, m_2, ..., m_p$ die Ziffern zur Basis $b$.
\item $e$ … \emph{Exponent}, ganzzahlig $e_{min}\leq e \leq e_{max}$.
\end{itemize}
In jedem Gleitkommasystem sind nur endlich viele Zahlen darstellbar, die Menge der reellen zahlen ist aber überabzählbar (unendlich).\\
Gleitkommazahlen liegen auf der Zahlengeraden diskret verteilt (fester Exponent $\curvearrowright$ gleiche Abstände, wächst Exponent um $k$, so wachsen die Abstände auf $b^k$-fache!)\\
Veranschaulichung für $b=10$, Mantissenlänge $p=1$:\\
\begin{tikzpicture}[scale=0.9]
\draw (0,0) -- (11,0);
\draw (0,0) node[below]{0};
\draw (0,0.15)--(0,-0.15);
\draw (1,0.15)--(1,-0.15);
\draw (2,0.15)--(2,-0.15);
\draw (3,0.15)--(3,-0.15);
\draw (4,0.15)--(4,-0.15);
\draw (5,0.15)--(5,-0.15);
\draw (6,0.15)--(6,-0.15);
\draw (7,0.15)--(7,-0.15);
\draw (8,0.15)--(8,-0.15);
\draw (9,0.15)--(9,-0.15);
\draw (10,0.15)--(10,-0.15);
\draw (1,0) node[below]{\textcolor{red}{1}};
\draw (2,0) node[below]{\textcolor{red}{2}};
\draw (3,0) node[below]{\textcolor{red}{3}};
\draw (4,0) node[below]{\textcolor{red}{4}};
\draw (5,0) node[below]{\textcolor{red}{5}};
\draw (6,0) node[below]{\textcolor{red}{6}};
\draw (7,0) node[below]{\textcolor{red}{7}};
\draw (8,0) node[below]{\textcolor{red}{8}};
\draw (9,0) node[below]{\textcolor{red}{9}};

\draw[green] (0.1,0.15)--(0.1,-0.15);
\draw[green] (0.2,0.15)--(0.2,-0.15);
\draw[green] (0.3,0.15)--(0.3,-0.15);
\draw[green] (0.4,0.15)--(0.4,-0.15);
\draw[green] (0.5,0.15)--(0.5,-0.15);
\draw[green] (0.6,0.15)--(0.6,-0.15);
\draw[green] (0.7,0.15)--(0.7,-0.15);
\draw[green] (0.8,0.15)--(0.8,-0.15);
\draw[green] (0.9,0.15)--(0.9,-0.15);


\draw (10,0) node[below]{\textcolor{cyan}{$1\cdot 10^1$...}};
\end{tikzpicture}\\
\emph{Exponent} 0: $1\cdot 10^{\textcolor{red}{0}}, 2\cdot 10^{\textcolor{red}{0}}, ... , 9\cdot 10^{\textcolor{red}{0}}$\\
\emph{Exponent} -1: $1\cdot 10^{\textcolor{green}{-1}}, 2\cdot 10^{\textcolor{green}{-1}}, ... , 9\cdot 10^{\textcolor{green}{-1}}$\\
\emph{Exponent} 1: $1\cdot 10^{\textcolor{cyan}{1}}, 2\cdot 10^{\textcolor{cyan}{1}}, ... , 9\cdot 10^{\textcolor{cyan}{1}}$\\
\emph{Rundung}: Zahlen, die nicht in diesem „Raster“ enthalten sind, werden auf dei nächstgelegene darstellbare Gleitkommazahl gerundet. Falls die Zahl genau in der Mitte zwischen zwei darstellbaren Zahlen liegt, wird auf die gerade Zahl gerundet (Bsp. $3,75 \rightarrow 3,8$ oder $4,65 \rightarrow 4,6$ bei Rundung auf eine Stelle nach dem Komma).
\paragraph{Numerische Probleme beim Rechnen mit Gleitkommazahlen}
\begin{itemize}
\item Kommutativ-, Assoziativ- und Distributivgesetze gelten im allgemeinen nicht mehr. Ursachen sind bspw. Ziffernauslöschung bei der Subtraktion fast gleicher Zahlen, Addition oder Subtraktion von Zahlen unterschiedlicher Größenordnung.
\subparagraph{Bsp. 11:}
\begin{enumerate}[label=\arabic*.)]
\item Man berechne $(a+b)+c$ und $a+(b+c)$ in einem System mit 3-Stelliger Mantisse:\\
$a=3,73\cdot 10^6$, $b=-3,71\cdot 10^6$ und $c=6,42\cdot 10^3$
\begin{itemize}
\item $a+b= 0,02\cdot 10^6=2,00 \cdot 10^4$ (Normalisierung)\\
$c=0,642\cdot 10^4=0,64\cdot 10^4$ (Exponentenangleichung und Rundung)\\
$(a+b)+c=2,64\cdot 10^4=\resultul{26400}$
\item $c=0,00642\cdot 10^6=0,01\cdot 10^6$ (Exponentenangleichung und Rundung)\\
$b+c=-3,,70\cdot 10^6$\\
$a+(b+c)=0,03\cdot 10^6=3,00\cdot 10^4 = \resultul{30000}$
\item exakter Wert: $a+b+c = \resultul{26420}$
\end{itemize}
\item Aufgabe der numerischen Mathematik ist es, die unvermeidlichen Genauigkeitsverluste beim Rechnen mit Maschinenzahlen durch optimale Organisation (Reihenfolge) der Rechnung und Fehleranalyse in Grenzen zu halten.
\end{enumerate}
\end{itemize}
\item Gleitkommaformat IEEE 754 (single precision, 32 Bit)\\
$x=v\cdot m \cdot b^e=(-1)^{\torange{V}}\cdot 1, \tblue{m_2\,m_3...m_24}\cdot 2^{\tred{E}-B}$ ($b=2$, Binärsystenm)
\begin{itemize}
\item Vorzeichen $V=0\curvearrowright$ positiv, $V=1\curvearrowright$ negativ (1 Bit)
\item Mantisse $m_1$ im Binärsystem stets gleich $1$.\\
$\curvearrowright$ nur Abspeicherung von $\tblue{\boxed{M=m_2\,m_3...m_24}}$ (23 Bit)
\item Exponent: Abgespeichert wird $\tred{\boxed{E:=e+B}} $ \\
mit dem sogenannten \emph{Biaswert $B=127$} (Bias = Verzerrung) als nichtnegative 8-stellige Binärzahl, $e_{min}=-126 $ ($E=1$), $e_{max}=127$ ($E=254=(1111\;1110)_2)$, die Grenzfälle $E=(0000\;0000)_2$ und $E=(1111\;1111)_2$ sind für Sonderfälle ($0$, $\infty$, nichtdefinierte Werte) vorgesehen.
\end{itemize}
Abspeicherung in der Reihenfolge $\boxed{\torange{V}\tred{E}\tblue{M}}$ (32 Bit)
\subparagraph{Bsp. 12:} Umwandlung einer Dezimalzahl in das IEEE 754-Format (32-Bit), $x=435,9$ (vgl. Bsp. 6)
\begin{enumerate}[label=\arabic*.)]
\item Konvertierung in Dualzahl (unter Verwendung von Bsp. 6/Hexadez.)\\
$x=1B3,E\overline{6})_{16}=(1\; 1011\;0011,1100\;0110\;0110 ...)_2$
\item Normalisierte Gleitkommadarstellung, Mantisse mit 23 Stellen nach dem Komma, Kommaverschiebung um 8 Stellen.\\
$\curvearrowright x = 1, \underbrace{1011\;0011\;1110\;0110\;0110\;011}_{\tblue{M}}(0\; 0110...))_2\cdot 2^8$ \quad (Abrundung!)
\item Exponent $e=8\curvearrowright E=e+B=8+127=135=\underbrace{(1000\;0111)_2}_{E}$
\item Vorzeichenbit $V=0$ (da $x$ positiv)
\end{enumerate}
$\curvearrowright x: \quad \boxed{0 \; \tred{1000\;0111}\;\tblue{1011\;0011\;1110\;0110\;0110\;011}}$
\paragraph{Bsp. 13:} IEEE 754$\rightarrow$ Dezimalzahl\\
$\boxed{1\;\tred{1000\;0011}\;\tblue{0111\; 1100 \; 0000\;0000\;0000 \; 000}}$
\begin{enumerate}[label=\arabic*.)]
\item $E=(\tred{1000\;00111})_2=131\curvearrowright e=E-B=131-127=\resultul{4}$
\item $V=1\curvearrowright$ negativ, normalisierte Mantisse $1,\tblue{M}$\\
$\curvearrowright x=-(1,\tblue{011111})_2\cdot 2^4=-(10111,11)_2$\\
$\curvearrowright x = -23,75$
\end{enumerate}
Bemerkung: 
\begin{enumerate}[label=\arabic*.)]
\item Neben dem single-Format gibt es in IEEE 754 das double-Format (54 Bit, V=1Bit, E=11Bit, M=52 Bit, B=1023) sowie das erweiterbare Format
\item Zahlbereiche single: $1,401\cdot 10^{-45}...3,403\cdot 10^{38}$, double: $4,941\cdot 10^{-324}...1,798\cdot 10^{308}$
\end{enumerate}
\end{enumerate}
\subsubsection{Ordnungsstruktur}
\begin{itemize}
\item Durch $\leqq$ (auch $\leq$) ist auf $\mathbb{R}$ eine vollständige Ordnungsrelation erklärt.
\item Verträglichkeit mit der algebraischen Struktur (für alle $x,y,z\in \mathbb{R}$):
\begin{enumerate}[label=(\arabic*)]
\item $x\leq y \quad \Rightarrow \quad x+z \leqq y+z$
\item $(x\leq y) \wedge (z \geq 0 )\quad \Rightarrow \quad x \cdot z \leq y \cdot z$\\
$(x\leq y) \wedge (z \leq 0 )\quad \Rightarrow \quad x \cdot z \geq y \cdot z$
\end{enumerate}
Für die strikte Ordnung $<$ gilt:\\
$\boxed{(x<y)\wedge (z<0) \quad \Rightarrow \quad x \cdot > y \cdot z}$
\end{itemize}
\paragraph{Def. 8:} \parskp
Sei $x$ eine reele Zahl. Dann heißt $|x|:=\begin{cases}
x \quad \text{ für } x\geq 0\\
-x \quad x <0
\end{cases}$ der (absolute) Betrag von $x$.\\
ABB 5\\
Vorzeichenfunktion $sgn(x):= \begin{cases}
1 \quad x>0\\
0 \quad x=0\\
-1 \quad x<0
\end{cases}$
\subparagraph{Diskussion:} Es gilt:
\begin{enumerate}
\item $|a-b|$ = „Abstand der Zahlen $a$ und $b$ auf der Zahlengeraden“\\
ABB 6\\
(speziell: $|a|$ = „Abstand von $a$ zum Ursprung“)
\item $a = |a| \cdot sgn (a)$
\item $|a|=|-a|, |ab|=|a|\cdot |b|, \left|\frac{a}{b}\right|=\frac{|a|}{|b|}$ (falls $b\not = 0$)
\item $|a \pm b| \leq |a| + |b|$ (\emph{Dreiecksungleichung}) für alle $a,b \in \mathbb{R}$
\end{enumerate}
\paragraph{Lösung von Ungleichung}
\subparagraph{Bsp. 14:} (Ungleichung mit Beträgen)\\
Gesucht sei die Lösungmenge $L$ der reellen Zahlen, die die Ungleichung $|x-1|<3+\frac{1}{2}x$ (*) erfüllen.
\begin{itemize}
\item kritische Stelle(n): Nullstellen des Terms innerhalb der Betragsstriche d.h. $x=1 \curvearrowright$ Fallunterscheidung\\
ABB 7
\item 1. Fall: $x-1<0$ d.h. $x<1$\\
in (*): \\
$\tred{-}(x-1)<3+\frac{1}{2}x \Leftrightarrow \\
-\frac{3}{2}x<2 \Leftrightarrow \\
\resultul{x>-\frac{4}{3}}$\\
$\curvearrowright L_1=\left\lbrace x|(x<1) \wedge x>\frac{4}{3}\right\rbrace = \left(-\frac{4}{3},1\right)$
\item 2. Fall $x-1 \geq 0$ d.h. $x\geq 1$\\
in (*):\\
$x-1< 3+ \frac{1}{2}x \Leftrightarrow\\
\frac{1}{2}x<4 \Leftrightarrow\\
\resultul{x<8}$\\
$\curvearrowright L_2 = \{ x| (x\geq 1) \wedge (x <8)\}=[1,8)$
\item $\Rightarrow L = L_1 \cup L_2 = \resultul{\left(-\frac{4}{3},8\right)}$
\end{itemize}
\subparagraph{Bsp. 15:} (Ungleichung mit gebrochen rationalen Termen)\\
$\frac{x}{x+1}<1$ (*)
\begin{itemize}
\item kritische Stelle(n): Nenner-Nullstellen, hier: $x=-1$.\\
ABB 7 (äquiv. mit -1)
\item 1. Fall: $x<-1$\\
in (*):\\
$\Leftrightarrow x>x+1 \Leftrightarrow 0>1$ (Widerspruch)\\
$L_1=\emptyset$.
\item 2. Fall: $x > -1$\\
in (*):\\
$\Leftrightarrow x < x+1 \Leftrightarrow 0<1$ (wahre Aussage)\\
$L_2=\{x|x>-1\}=(-1,\infty)$
\item $\Rightarrow L = L_1 \cup L_2 = \resultul{(-1,\infty)}$
\end{itemize}
\subparagraph{Bsp. 16:} (quadratische Ungleichungen)\\
$x^2+3x<10 \Leftrightarrow \\
\left( x+\frac{3}{3}\right)^2-\frac{9}{4}<10 \Leftrightarrow$ (vereinfacht durch quadratische Ergännzung)\\
$\left( x+\frac{3}{2}\right)^2<\frac{49}{4} \Leftrightarrow\\
\left|x+\frac{3}{2}\right|<\frac{7}{2} \Leftrightarrow$ (Äquivalenz siehe Übung) \\
$\frac{-7}{2}<x+\frac{3}{2}<\frac{7}{2} \Leftrightarrow\\
-5<x<2$\\
$L=\resultul{(-5,2)}$


\end{document}